\documentclass[UTF8,12pt,a4paper]{ctexart}

\usepackage[a4paper,margin=2.5cm]{geometry}
\usepackage{amsmath,amssymb,bm}
\usepackage{booktabs,longtable,tabularx,array,multirow}
\usepackage{graphicx,float,caption,subcaption}
\usepackage{siunitx}
\usepackage{enumitem}
\usepackage{xcolor}
\usepackage{fancyhdr}
\usepackage{fvextra}
\usepackage{listings}
\usepackage{hyperref}

\hypersetup{hidelinks}
\urlstyle{same}
\IfFontExistsTF{Microsoft YaHei}
  {\setCJKmonofont{Microsoft YaHei}}
  {\setCJKmonofont{FandolFang-Regular}}
\IfFontExistsTF{JuliaMono}
  {\setmonofont{JuliaMono}[Scale=0.875]}
  {\IfFontExistsTF{Noto Sans Mono}
    {\setmonofont{Noto Sans Mono}[Scale=MatchLowercase]}
    {\IfFontExistsTF{DejaVu Sans Mono}{\setmonofont{DejaVu Sans Mono}[Scale=MatchLowercase]}{}}}
\definecolor{codeblue}{RGB}{31,78,121}
\definecolor{codegreen}{RGB}{40,110,85}
\definecolor{codegray}{RGB}{95,105,115}
\lstdefinestyle{modelpython}{
  language=Python,
  numbers=left,
  numberstyle=\tiny\color{codegray},
  numbersep=8pt,
  basicstyle=\ttfamily\fontsize{8.0}{9.2}\selectfont,
  keywordstyle=\color{codeblue}\bfseries,
  commentstyle=\color{codegreen},
  stringstyle=\color{purple!70!black},
  breaklines=true,
  breakatwhitespace=false,
  frame=single,
  rulecolor=\color{black!38},
  columns=fullflexible,
  keepspaces=true,
  linewidth=\textwidth,
  showstringspaces=false,
  tabsize=4,
  xleftmargin=2.3em,
  framexleftmargin=0pt,
  aboveskip=3pt,
  belowskip=3pt
}
\graphicspath{{../figures/}}
\captionsetup{font=small,labelsep=quad}
\setlength{\parindent}{2em}
\setlength{\parskip}{0pt}
\setlist{nosep}
\pagestyle{plain}
\newcommand{\TODO}[1]{\textcolor{red}{\textbf{TODO[#1]}}}
\newcommand{\papertitle}{TODO[填写能够概括方法与对象的论文标题]}

\begin{document}

\begin{center}
  {\zihao{2}\bfseries \papertitle}
\end{center}

\vspace{1.2em}
\noindent\textbf{摘要：}
\TODO{写清研究对象、每问采用的方法、关键数值结果及主要检验；不写背景套话。}

\vspace{0.8em}
\noindent\textbf{关键词：}
\TODO{填写3--5个能够检索模型与研究对象的关键词}

\clearpage
\setcounter{page}{2}

\section{问题重述}

\TODO{用自己的语言凝练任务、对象、输入与交付物，不逐句抄写题面。}

\section{问题分析}

\TODO{说明各小题之间的数据流、约束关系与总体求解路线。若依赖关系复杂，插入一张信息密度足够的流程图。}

\section{模型假设}

\begin{enumerate}
  \item \TODO{给出必要假设，并说明它会排除什么情形或带来何种偏差。}
\end{enumerate}

\section{符号说明与数据预处理}

\subsection{主要符号}

\begin{table}[H]
  \centering
  \caption{主要符号及含义}
  \begin{tabular}{cll}
    \toprule
    符号 & 含义 & 单位 \\
    \midrule
    \TODO{符号} & \TODO{含义} & \TODO{单位} \\
    \bottomrule
  \end{tabular}
\end{table}

\subsection{数据审计与预处理}

\TODO{报告数据规模、缺失、异常、重复、单位转换和预处理原则，并保留关键诊断证据。}

% ORBIT:QUESTIONS:BEGIN
% 问题分节由 bootstrap_project.py 在此处增补；完成后直接在本文件中写作。
% ORBIT:QUESTIONS:END

\section{模型评价与推广}

\subsection{模型优点}

\TODO{结合本题证据说明可解释性、精度、稳健性或计算效率，避免空泛形容。}

\subsection{模型局限与改进}

\TODO{指出假设、数据、算法或外推范围造成的具体限制，并给出可执行改进方向。}

\section*{AI工具使用声明}

本参赛队在竞赛过程中使用了AI工具，主要用于\TODO{简要填写实际用途，如语言润色、代码调试等}，详细使用情况见支撑材料。

{\footnotesize
\begin{thebibliography}{99}
\setlength{\itemsep}{0.2em}

\bibitem{todo-reference}
\TODO{按正文首次引用顺序补充真实、可核验的参考文献。}

\end{thebibliography}
}
\label{body-end}

\clearpage
\appendix
\IfFileExists{../.cumcm/generated/support-files.tex}{\begin{itemize}
  \item \TODO{运行构建脚本后自动生成支撑材料列表}
\end{itemize}
}{}
\IfFileExists{../.cumcm/generated/source-code.tex}{\TODO{运行构建脚本后自动生成代码附录}
}{}

\end{document}
